%! Rational Functions and End Behavior — Unit 3, Lesson 3.1 — Lesson Plan
\documentclass[10pt]{article}
\usepackage{precalculus-article}
\usepackage{precalculus-boxes}
\usepackage{graphicx}
\graphicspath{{images/}}
\newcommand{\TallMath}[1]{$\displaystyle #1\rule[-1.4em]{0pt}{3.2em}$}
\newcommand{\CourseName}{Precalculus}
\newcommand{\SchoolYear}{2026--2027}
\newcommand{\MeetingLength}{55 minutes}
\newcommand{\UnitNumberName}{Unit 3: Rational Functions \quad}
\newcommand{\LessonNumberName}{Lesson 3.1: Rational Functions and End Behavior}

\begin{document}
\begin{center}
    {\Large\bfseries \CourseName: \SchoolYear \\
    {\small\bfseries \UnitNumberName \LessonNumberName}}
\end{center}

\begin{tcolorbox}[colback=sky, colframe=navy, boxrule=0.9pt, arc=2mm,
    left=2mm, right=2mm, top=1.5mm, bottom=1.5mm]
    \textbf{Primary Objective:} Students will describe the end behavior of rational functions
    by examining the quotient of leading terms, identify horizontal and slant asymptotes,
    and express end behavior using limit notation.
    \hfill\textit{\small Skills: MP 1.B, 3.A}
\end{tcolorbox}

\renewcommand{\arraystretch}{1.6}
\begin{skillbox}[Priority Ideas \& Skills]{goldbox}
\begin{tabularx}{\linewidth}{@{}X X@{}}
\textbf{AP Skills} &
\textbf{Key Understandings} \\
\midrule
\textbf{MP 1.B} --- Identify mathematical information from graphical, numerical,
analytical, and verbal representations. &
\begin{itemize}[leftmargin=1.2em,itemsep=2pt,topsep=0pt]
  \item A rational function $r(x) = p(x)/q(x)$ measures the relative size of
        two polynomials. (EK 1.7.A.1)
  \item For large $|x|$, leading terms dominate; end behavior equals the quotient
        of leading terms. (EK 1.7.A.2)
  \item deg num $>$ deg denom $\Rightarrow$ polynomial end behavior; if the
        difference in degrees is 1, a slant asymptote exists. (EK 1.7.A.3)
\end{itemize} \\
\textbf{MP 3.A} --- Describe characteristics of a function using appropriate
mathematical language and notation. &
\begin{itemize}[leftmargin=1.2em,itemsep=2pt,topsep=0pt]
  \item deg num $=$ deg denom $\Rightarrow$ HA at $y = a_n / b_n$. (EK 1.7.A.4)
  \item deg num $<$ deg denom $\Rightarrow$ HA at $y = 0$. (EK 1.7.A.5)
  \item HAs are described with limit notation:
        $\lim_{x\to\infty} r(x) = b$. (EK 1.7.A.6)
\end{itemize} \\
\end{tabularx}
\end{skillbox}

\begin{skillbox}[{Vocabulary, Concepts, \& Theorems}]{greenbox}
\renewcommand{\arraystretch}{1.4}
\begin{tabularx}{\linewidth}{@{} >{\bfseries\color{navy}}p{0.26\linewidth} X @{}}
rational function         & A function of the form $r(x) = p(x)/q(x)$ where $p$ and $q$ are polynomials. \\
quotient of leading terms & The ratio of the leading terms of numerator and denominator; governs end behavior. \\
horizontal asymptote      & A line $y = b$ that the graph approaches (and stays near) as $x \to \pm\infty$. \\
slant asymptote           & A non-horizontal line that the graph approaches when deg num $=$ deg denom $+ 1$. \\
end behavior              & The output trend as $x \to +\infty$ or $x \to -\infty$. \\
limit notation            & $\lim_{x\to\infty} r(x) = b$: outputs get arbitrarily close to $b$ as $x$ grows. \\
\end{tabularx}
\end{skillbox}

\begin{skillbox}[Activate Prior Knowledge \& Spiral Review]{sky}
    \begin{tabularx}{\linewidth}{@{}X c@{}}
        \begin{minipage}[t]{0.70\linewidth}\vspace{0pt}
            Please see attached warm-up (spiral review from Lesson 1.6).
            Skills reviewed: end behavior of polynomial functions using the leading term,
            limit notation for polynomials, and degree/leading-coefficient analysis.
        \end{minipage}
        &
        \begin{minipage}[t]{0.24\linewidth}\vspace{0pt}\centering
            \textit{\small (authored warm-up compiles to target/)}
        \end{minipage}
    \end{tabularx}
\end{skillbox}

\begin{skillbox}[Hook]{sky}
    A therapist charges \$200 per session plus a one-time \$50 intake fee.
    The \textbf{average cost per session} after $n$ sessions is:
    \[
      r(n) = \frac{200n + 50}{n}
    \]
    \textit{Discussion prompts:}
    \begin{enumerate}[itemsep=2pt]
        \item Evaluate $r(1)$, $r(10)$, $r(100)$. What pattern do you notice?
        \item As you attend more and more sessions, what does the average cost approach? Why?
        \item Can the average cost ever equal exactly \$200? What does that tell us about the graph?
    \end{enumerate}
    \textit{Takeaway:} The \$50 intake fee is ``spread out'' over more sessions, so the average
    approaches \$200 but never reaches it --- the graph has a \textbf{horizontal asymptote} at $y = 200$.
    This is the quotient of leading terms: $200n/n = 200$.
\end{skillbox}

\begin{skillbox}[Lesson]{sky}
\begin{multicols}{2}

\textbf{Part 1: Quotient of Leading Terms} (10 min)

Key insight: for large $|x|$, lower-degree terms become negligible compared to the
leading term of each polynomial.  So end behavior is controlled by the
\textit{quotient of the leading terms}.

\medskip
\textbf{Worked Example.}
\[
  r(x) = \frac{3x^2 + 1}{x^2 - 4}
\]
\begin{itemize}
  \item Leading term of numerator: $3x^2$
  \item Leading term of denominator: $x^2$
  \item Quotient: $\dfrac{3x^2}{x^2} = 3$
  \item As $x \to \pm\infty$, $r(x) \to 3$.
  \item Horizontal asymptote: $y = 3$.
\end{itemize}

Numerical check: $r(100) = \dfrac{30{,}001}{9{,}996} \approx 3.001$.

\columnbreak

\textbf{Part 2: The Three Cases} (10 min)

\renewcommand{\arraystretch}{1.5}
\begin{tabular}{@{}p{2.4cm} p{3.6cm}@{}}
\textbf{Degree comparison} & \textbf{End behavior} \\
\hline
deg num $<$ deg denom & HA at $y = 0$ \\
deg num $=$ deg denom & HA at $y = a_n/b_n$ \\
deg num $>$ deg denom & Polynomial (slant if diff.\,$= 1$) \\
\end{tabular}

\medskip
Quick examples:
\begin{itemize}
  \item $\dfrac{x+1}{x^2-9}$: $1 < 2$; HA at $y = 0$.
  \item $\dfrac{3x^2}{x^2+4}$: $2 = 2$; HA at $y = 3$.
  \item $\dfrac{x^3-2x}{x-3}$: $3 > 1$; polynomial ($\approx x^2$) end behavior.
\end{itemize}

\medskip
\textbf{Part 3: Limit Notation for HAs} (10 min)

When $r$ has a HA at $y = b$:
\[
  \lim_{x\to\infty} r(x) = b \qquad \lim_{x\to -\infty} r(x) = b
\]
Example: for $r(x) = \dfrac{3x^2+1}{x^2-4}$,
\[
  \lim_{x\to\infty} r(x) = 3 \quad \text{and} \quad \lim_{x\to -\infty} r(x) = 3.
\]
For $\dfrac{x+1}{x^2-9}$: both limits equal $0$.

\end{multicols}
\end{skillbox}

\begin{skillbox}[Active Monitoring]{redbox}
While students work on guided notes and the activity, circulate and check:
\begin{itemize}
  \item \textbf{Common error:} Students look at constant terms (e.g., the ``$+1$'' or ``$-4$'')
        instead of leading terms.  Ask: ``Which term dominates for very large $x$?''
  \item \textbf{Common error:} Confusing a horizontal asymptote with a zero of the denominator.
        Clarify: HAs come from degree comparison, not from roots of $q(x)$.
  \item \textbf{Notation check:} Are students writing $\lim_{x\to\infty}$ correctly?
  \item \textbf{Cold-call:} ``If the numerator has degree 4 and the denominator degree 6,
        what is the horizontal asymptote?'' ($y = 0$.)
  \item \textbf{Cold-call:} ``Give me a rational function with a HA at $y = -2$.''
\end{itemize}
\end{skillbox}

\begin{skillbox}[Group Work \& Differentiation]{redbox}
Students work on the \textbf{Group Activity: Analyzing Concert Venue Costs} in leveled tiers.
All tiers use the same four rational functions in context.

\begin{itemize}
  \item \textbf{Tier R (Remediate):} Leading terms are provided; students compute the quotient
        and state whether a HA exists (and its value).
  \item \textbf{Tier A (Apply):} Students identify degrees and leading coefficients from the
        full expressions, compute the quotient, write both limit statements, and name the asymptote type.
  \item \textbf{Tier E (Extend):} Students construct rational functions meeting specified HA or
        slant asymptote conditions, then verify with graphing technology.
\end{itemize}
\end{skillbox}

\begin{skillbox}[Individual Work \& Assessment]{redbox}
\textbf{Exit Ticket} (last 5--7 minutes, individual, collected):
\begin{enumerate}
  \item For $r(x) = \dfrac{5x^2 - 3}{2x^2 + x - 1}$: state the quotient of leading terms,
        write $\lim_{x\to\infty} r(x)$ and $\lim_{x\to -\infty} r(x)$, and identify the HA.
  \item A rational function has numerator degree 3 and denominator degree 5.
        State the HA and write the limit statement.
\end{enumerate}
\textbf{Assessment note:} Look for correct identification of \textit{leading} terms (not
constants), proper limit notation, and correct application of the three degree-comparison cases.
\end{skillbox}

\begin{skillbox}[Reinforcement \& Extension]{goldbox}
\textbf{Homework:} 5 problems covering limit notation for four rational functions, conceptual
explanation of why equal-degree quotients approach $a_n/b_n$, interpreting a HA at $y = 0$,
describing end behavior with a slant asymptote, and one AP-style multiple-choice item.

\textbf{Extension (optional):} Students construct and verify their own rational functions with
prescribed asymptotic behavior, using graphing technology to confirm predictions.

\textbf{Preview:} Lesson 1.8 examines \textit{zeros} of rational functions --- the input values
where the numerator equals zero --- and how those points interact with the asymptotes from today.
\end{skillbox}
\end{document}
